
\section{Experiments}

\subsection{Experimental Protocol}
The current implemented model is Prithvi-EO-2.0 300M TL. The backbone is frozen and only the decoder/head is trained. Additional candidate RSFMs to be evaluated under the same protocol are SatMAE++, SoftCon, DOFA, and SpectralGPT.

\begin{table}[t]
\centering
\caption{Prithvi-EO-2.0 linear-probing setup used in the current experiments.}
\label{tab:training_config}
\resizebox{\linewidth}{!}{
\begin{tabular}{l l}
\toprule
Parameter & Value \\
\midrule
Backbone & Prithvi-EO-2.0 300M TL \\
Trainable part & MLP decoder/head only \\
Frozen parameters & Encoder backbone \\
Bands & Blue, Green, Red, NIR Narrow, SWIR1, SWIR2 \\
Input size & $224\times224$ \\
Batch size & 128 \\
Optimizer & AdamW \\
Learning rate & $1\times10^{-3}$ with cosine annealing \\
Weight decay & $1\times10^{-4}$ \\
Precision & BF16 mixed precision \\
Early stopping & Patience = 15 epochs \\
\bottomrule
\end{tabular}}
\end{table}

\subsection{Current Global Results}
Table~\ref{tab:global_results} summarizes the current Prithvi-EO-2.0 results. The key observation is that strong global performance can coexist with large biome-wise disparity.

\begin{table}[t]
\centering
\caption{Global test performance of Prithvi-EO-2.0 linear probing. Values without completed experiments are placeholders.}
\label{tab:global_results}
\resizebox{\linewidth}{!}{
\begin{tabular}{l l l c c}
\toprule
Dataset & Task & Training & Metric & Value \\
\midrule
m-EuroSAT & Classification & Original & Macro-F1 & 0.936 \\
m-BigEarthNet & Multi-label cls. & Original & Macro-F1 & 0.626 \\
m-BigEarthNet & Multi-label cls. & Balanced & Macro-F1 & 0.424 \\
m-SA-Crop-Type & Segmentation & Original & Global mIoU & 0.2929 \\
m-SA-Crop-Type & Segmentation & Original & Pixel Acc. & 0.6512 \\
m-So2Sat & Classification & Original & Macro-F1 & \redtodo{TBD} \\
m-Cashew-Plantation & Segmentation & Original & mIoU & \redtodo{TBD} \\
Sen1Floods11 & Segmentation & Original & mIoU/F1 & \redtodo{TBD} \\
Multi-Crop-Seg. & Segmentation & Original & mIoU & \redtodo{TBD} \\
MMEarth100K & Segmentation & Original & mIoU & \redtodo{TBD} \\
\bottomrule
\end{tabular}}
\end{table}

\subsection{Per-Biome Results: Classification}
Table~\ref{tab:per_biome_classification} reports current per-biome results for m-EuroSAT and m-BigEarthNet. For m-EuroSAT, the model performs strongly on dominant temperate and Mediterranean biomes but drops on Boreal Forests/Taiga and Temperate Conifer Forests. For m-BigEarthNet, the Mediterranean biome has the lowest F1@opt under original training.

\begin{table*}[t]
\centering
\caption{Current per-biome classification results for Prithvi-EO-2.0 linear probing. F1@opt denotes macro-F1 at per-class optimal thresholds for the multi-label m-BigEarthNet task. A dash indicates that the experiment has not yet been run.}
\label{tab:per_biome_classification}
\resizebox{0.8\textwidth}{!}{
\begin{tabular}{l l c c c}
\toprule
Dataset & Biome & Test $N$ & Original & Balanced \\
\midrule
\multirow{6}{*}{m-EuroSAT Macro-F1}
& Temperate Broadleaf \& Mixed Forests & 706 & 0.955 & -- \\
& Mediterranean Forests, Woodlands \& Scrub & 231 & 0.948 & -- \\
& Boreal Forests/Taiga & 18 & 0.800 & -- \\
& Temperate Conifer Forests & 25 & 0.818 & -- \\
& Temperate Grasslands, Savannas \& Shrublands & 8 & 1.000 & -- \\
& Unknown / water & 12 & 0.871 & -- \\
\midrule
\multirow{5}{*}{m-BigEarthNet F1@opt}
& Temperate Broadleaf \& Mixed Forests & 403 & 0.621 & 0.453 \\
& Mediterranean Forests, Woodlands \& Scrub & 336 & 0.492 & 0.402 \\
& Boreal Forests/Taiga & 132 & 0.668 & 0.525 \\
& Temperate Conifer Forests & 14 & 0.625 & 0.605 \\
& Unknown / water & 115 & 0.551 & -- \\
\bottomrule
\end{tabular}}
\end{table*}

\subsection{Per-Biome Results: Segmentation}
Table~\ref{tab:per_biome_segmentation} reports current segmentation results for m-SA-Crop-Type. Performance is substantially lower in Deserts and Xeric Shrublands than in Mediterranean Forests, Woodlands and Scrub.

\begin{table}[t]
\centering
\caption{Current per-biome segmentation results for Prithvi-EO-2.0 on m-SA-Crop-Type.}
\label{tab:per_biome_segmentation}
\resizebox{\linewidth}{!}{
\begin{tabular}{l c c}
\toprule
Biome & Test $N$ & mIoU \\
\midrule
Deserts \& Xeric Shrublands & 124 & 0.1755 \\
Mediterranean Forests, Woodlands \& Scrub & 875 & 0.2891 \\
\bottomrule
\end{tabular}}
\end{table}

\subsection{Fairness Metrics Across Biomes}
Table~\ref{tab:fairness_summary} summarizes biome fairness metrics. Lower NFR, standard deviation, EOdd, and DPM indicate more consistent behavior across biomes. The balanced m-BigEarthNet run reduces neither the worst-biome metric nor the NFR, suggesting that naive biome balancing is not sufficient.

\begin{table*}[t]
\centering
\caption{Current biome fairness summary for Prithvi-EO-2.0 linear probing. Worst-metric is worst-biome macro-F1 for classification and worst-biome mIoU for segmentation.}
\label{tab:fairness_summary}
\resizebox{0.9\textwidth}{!}{
\begin{tabular}{l l c c c c c}
\toprule
Dataset & Training & Worst-metric & NFR $\downarrow$ & Std. dev. $\downarrow$ & EOdd $\downarrow$ & DPM $\downarrow$ \\
\midrule
m-EuroSAT & Original & 0.800 & 0.2212 & 0.080 & 0.247 & 0.234 \\
m-BigEarthNet & Original & 0.492 & 0.338 & 0.077 & 0.429 & 0.168 \\
m-BigEarthNet & Balanced & 0.402 & 0.463 & 0.079 & 0.407 & 0.166 \\
m-SA-Crop-Type & Original & 0.1755 & 0.4891 & 0.0568 & -- & -- \\
m-So2Sat & Original & \redtodo{TBD} & \redtodo{TBD} & \redtodo{TBD} & \redtodo{TBD} & \redtodo{TBD} \\
m-Cashew-Plantation & Original & \redtodo{TBD} & \redtodo{TBD} & \redtodo{TBD} & -- & -- \\
\bottomrule
\end{tabular}}
\end{table*}

\subsection{Biome Distribution in m-BigEarthNet}
The current m-BigEarthNet split is highly imbalanced across biomes. Temperate Broadleaf and Mixed Forests and Mediterranean regions dominate the original test split, while Temperate Conifer Forests contain only 1.4\% of test samples. Table~\ref{tab:ben_distribution} shows the current distribution.

\begin{table}[t]
\centering
\caption{Biome distribution in the current m-BigEarthNet preprocessing.}
\label{tab:ben_distribution}
\resizebox{\linewidth}{!}{
\begin{tabular}{l c c c c}
\toprule
Biome & Train bal. & Val orig. & Val bal. & Test orig. \\
\midrule
Boreal Forests/Taiga & 356 & 139 & 17 & 132 \\
Mediterranean & 356 & 335 & 17 & 336 \\
Temp. Broadleaf \& Mixed & 356 & 379 & 17 & 403 \\
Temp. Conifer & 356 & 17 & 17 & 14 \\
Unknown & -- & 130 & -- & 115 \\
\midrule
Total & 1424 & 1000 & 68 & 1000 \\
\bottomrule
\end{tabular}}
\end{table}

\subsection{Planned Full Model Benchmark}
Table~\ref{tab:model_placeholder} gives the final model-comparison format. Current results are available for Prithvi-EO-2.0; other RSFMs will be filled after running the unified linear-probing pipeline.

\begin{table*}[t]
\centering
\caption{Planned full RSFM comparison on \method. Fill the placeholders after running the unified protocol.}
\label{tab:model_placeholder}
\resizebox{\textwidth}{!}{
\begin{tabular}{l ccc ccc}
\toprule
\multirow{2}{*}{Model} & \multicolumn{3}{c}{Classification Avg. / Worst / NFR} & \multicolumn{3}{c}{Segmentation Avg. / Worst / NFR} \\
\cmidrule(lr){2-4}\cmidrule(lr){5-7}
& Avg. & Worst & NFR $\downarrow$ & Avg. & Worst & NFR $\downarrow$ \\
\midrule
Prithvi-EO-2.0 & \redtodo{aggregate} & \redtodo{aggregate} & \redtodo{aggregate} & \redtodo{aggregate} & \redtodo{aggregate} & \redtodo{aggregate} \\
SatMAE++ & \redtodo{TBD} & \redtodo{TBD} & \redtodo{TBD} & \redtodo{TBD} & \redtodo{TBD} & \redtodo{TBD} \\
SoftCon & \redtodo{TBD} & \redtodo{TBD} & \redtodo{TBD} & \redtodo{TBD} & \redtodo{TBD} & \redtodo{TBD} \\
DOFA & \redtodo{TBD} & \redtodo{TBD} & \redtodo{TBD} & \redtodo{TBD} & \redtodo{TBD} & \redtodo{TBD} \\
SpectralGPT & \redtodo{TBD} & \redtodo{TBD} & \redtodo{TBD} & \redtodo{TBD} & \redtodo{TBD} & \redtodo{TBD} \\
\bottomrule
\end{tabular}}
\end{table*}
